\chapter{Elementary foundations and exact completion}\label{ch:foundations}
\sourcebox{Original olympiad prompt; first helper note; later proof notes}{The starting powerful-core argument was supplied by the user with attribution to Nikola Gyulev. The displayed multiplicativity line in the original prompt had a typographical duplication; the corrected identity is used here.}

\section{Local identities and safe cancellation}
For a prime $p$ and $e\ge0$,
\[
 \sigma(p^e)=\frac{p^{e+1}-1}{p-1},\qquad
 h(p^e)=p^e\frac{p^{e+1}-1}{p-1},\qquad
 p\nmid\sigma(p^e).
\]
Both $\sigma$ and $h$ are multiplicative on coprime arguments. Every preimage
$k$ divides $N$ and satisfies $k^2\le N$, strictly for $k>1$. We have $f(1)=1$.
If $h(a)=h(b)$, the product
\[
 C=\prod_{v_p(a)=v_p(b)>0}p^{v_p(a)}
\]
is unitary in both inputs, so $h(a/C)=h(b/C)$. Dividing the inputs by their
arbitrary gcd and canceling its $h$-value is generally invalid.

\section{Squarefree injectivity and the powerful core}
\begin{lemma}\label{found:squarefree}
If $a,b$ are squarefree and $h(a)=h(b)$, then $a=b$.
\end{lemma}
\begin{proof}
Cancel all common prime blocks $p(p+1)$. The residual prime supports are disjoint.
If the largest remaining prime is $q\ge5$, assume it occurs on the first side.
Every opposite prime $p$ has $p<q$ and $p+1<q$: equality would make the even
integer $q-1\ge4$ prime. Thus $q$ cannot divide the opposite product, a
contradiction. The only remaining bases could be $2,3$, whose squarefree
products have distinct $h$-values $1,6,12,72$. The residuals must be equal.
\end{proof}
This is also a specialization of Noppakaew--Pongsriiam's Theorem 12 \cite{np}.

Write
\[
 d(k)=\prod_{v_p(k)\ge2}p^{v_p(k)},\qquad s(k)=k/d(k).
\]
The second factor is squarefree and coprime to the first. Hence
$N=h(d(k))h(s(k))$. A fixed $d(k)$ permits at most one $s(k)$ by
Lemma~\ref{found:squarefree}. At a target prime $p$, the core exponent belongs
to $\{0,2,3,\ldots,\alpha_p\}$, a set of $\alpha_p$ choices. Therefore
\begin{equation}\label{found:product}
 \boxed{f(N)\le\prod_{p\mid N}\alpha_p.}
\end{equation}
The exact statement retains the completion condition:
\begin{equation}\label{found:completion}
 f(N)=\#\left\{d:\begin{array}{l}
 d\text{ powerful},\ h(d)\mid N,\\
 N/h(d)=h(s)\text{ for squarefree }s,\ \gcd(d,s)=1
 \end{array}\right\}.
\end{equation}
In particular, an upper bound on arbitrary possible cores can be much too large.

\section{The original maximal-order constant}
Let $g(N)=\prod_{p\mid N}\alpha_p$ and $c=\log3/3$. For every integer $a\ge1$,
$\log a\le ca$. Set $L=\log N$, $w=\log L$, and $z=L/w^3$. The primes at most
$z$ contribute at most $O(z w)=O(L/w^2)$ to $\log g(N)$. For the others,
\[
 \sum_{p>z}\log\alpha_p\le c\sum_{p>z}\alpha_p\le cL/\log z.
\]
Thus $\log f(N)\le(c+o(1))L/w$. Along
$N_y=\prod_{p\le y}p^3$, the prime number theorem gives
$\log g(N_y)\sim c\log N_y/\log\log N_y$.
This establishes sharpness for the majorant $g$, not for the true $f$.
Kominers gives the corresponding direct theorem and discussion \cite{kominers}.

\section{A deterministic squarefree inverse}
A squarefree inverse of a positive integer residual $R$, if it exists, is unique.
If $P^+(R)=q\ge5$, any squarefree preimage must contain $q$: a smaller base
cannot generate $q$ through its factor $p(p+1)$, while a larger base would
itself divide $R$. Divide $R$ by $q(q+1)$, rejecting a nonintegral division,
and record $q$. Reject a repeated base. Repeat until no prime at least $5$
remains. The terminal residual must be one of
\[
 (s,h(s))=(1,1),(2,6),(3,12),(6,72).
\]
Every accepted step is forced, and the final product verifies the inverse.
When the procedure is used in \eqref{found:completion}, also enforce
$\gcd(d,s)=1$.

\section{Integer coordinates adapted to the squarefree part}
The numbers $h(p)=p(p+1)$, indexed by primes, form a multiplicative
$\mathbb Z$-basis of the positive rational numbers. Namely, every $x\in\mathbb Q_{>0}$
has a unique finite expression
\[
 x=\prod_p h(p)^{b_p(x)},\qquad b_p(x)\in\mathbb Z.
\]
For $p\ge5$, all prime factors of $p+1$ are smaller than $p$, so descending
elimination is triangular with diagonal coefficient one. The initial two
prime coordinates are invertible over $\mathbb Z$ because
\[
 2=\frac{h(3)}{h(2)},\qquad 3=\frac{h(2)^2}{h(3)}.
\]
A squarefree inverse exists exactly when every $b_p(x)$ belongs to $\{0,1\}$.
Coprimality with a chosen core adds $b_p(x)=0$ for every $p\mid d$.
For example,
\[
 h(4)=\frac{h(2)h(7)}{h(3)}
\]
encodes the collision $h(12)=h(14)$. This coordinate change gives an exact
completion constraint; no asymptotic advantage has been proved from it alone.

\section{Why counting divisible cores cannot finish the problem}\label{found:badcores}
Define $B(N)=\#\{d:d\text{ powerful},\ h(d)\mid N\}$. The desired little-$o$
bound is false for $B$. For
\[
 N_y=\prod_{3\le p\le y}h(p^2)=h\!\left(\prod_{3\le p\le y}p^2\right),
\]
every subset $T$ gives a distinct core $d_T=\prod_{p\in T}p^2$ with
$h(d_T)\mid N_y$. Thus $B(N_y)\ge2^{\pi(y)-1}$. Since
$\log N_y=(4+o(1))y$,
\[
 \log B(N_y)\ge\left(\frac{\log2}{4}+o(1)\right)
 \frac{\log N_y}{\log\log N_y}.
\]
These targets are odd and represented. Among these particular subset cores,
only the full subset admits squarefree completion: a proper subset leaves an
odd residual greater than one, whereas $h(s)$ is even for every squarefree
$s>1$. This refutes the relaxed count, not the original conjecture.

\section{Lower-bound mechanisms that every proof must permit}
If $h(a)=h(b)=M$ and $\gcd(d,ab)=1$, then $h(da)=h(db)=h(d)M$.
This propagates one collision but does not create independent choices.
If $h(a_i)=h(b_i)=M_i$, $a_i\ne b_i$, and the prime supports of $a_ib_i$
are pairwise disjoint, choosing either member of every pair gives $2^r$
distinct preimages of $\prod_{i=1}^rM_i$.
An infinite construction with
\[
 \log\prod_{i=1}^rM_i=O(r\log r)
\]
would refute the primary conjecture. No such family was obtained.

For distinct Mersenne primes $M_i=2^{r_i}-1$, put
\[
 R=\sum_i r_i,\qquad k_i=2^{r_i-1}\prod_{j\ne i}M_j.
\]
Then
\[
 \sigma(k_i)=M_i2^{R-r_i},\qquad
 \boxed{h(k_i)=2^{R-1}\prod_jM_j.}
\]
The $k_i$ are distinct. This is the corrected form of the Mersenne construction
used throughout; it must not be quoted without the product over $j\ne i$.
No infinitude assumption about Mersenne primes is used.
